% =====================================================================
\chapter*{Instrucciones}
\addcontentsline{toc}{chapter}{Instrucciones}
\markboth{Instrucciones}{}
% =====================================================================

Este examen evalúa la fase~VI y cierra el curso. Cubre la formulación y resolución de problemas de optimización (módulo~24), la optimización bayesiana y los modelos sustitutos (módulo~25), el control predictivo y el aprendizaje por refuerzo (módulo~26), la optimización geométrica y 3D (módulo~27) y la búsqueda con modelos de lenguaje, la puesta en producción y la cuantificación del valor (módulo~28). El examen escrito (Partes~A y~B) comprueba el dominio de las herramientas a mano; el proyecto final (Parte~C), que es el núcleo del hito, comprueba que el estudiante sabe diseñar, construir y defender un optimizador con IA completo.

\section*{Condiciones}
\begin{center}
\small
\begin{tabularx}{\textwidth}{@{}lX@{}}
\toprule
Duración & 3 horas (Partes~A y~B). \\
Material permitido & Calculadora científica no programable y una hoja A4 de fórmulas escrita a mano por ambas caras. No se permiten apuntes, libros ni dispositivos conectados. Las tablas y los datos de referencia de estas instrucciones forman parte del examen. \\
Puntuación del examen & 100 puntos: Parte~A, 40 puntos (20 cuestiones de 2 puntos, 4 por módulo); Parte~B, 60 puntos (5 problemas de 12 puntos, uno por módulo). \\
Aprobado del examen & 70 puntos (70\,\%) o más. \\
Proyecto final (Parte~C) & Se califica aparte sobre 100 puntos con la rúbrica de la opción elegida (\cref{cap:proyecto}); dura 8 semanas y se aprueba con 70 puntos o más y todos los criterios de aceptación obligatorios cumplidos. \\
Calificación del hito & Nota del hito $=0{,}4\cdot\text{examen}/10+0{,}6\cdot\text{proyecto}/10$, con al menos un 7 en cada una de las dos partes. \\
Módulos & M24, M25, M26, M27 y M28 (reparto por módulo en el \cref{ap:autoevaluacion}). \\
\bottomrule
\end{tabularx}
\end{center}

\section*{Normas}
\begin{enumerate}[label=\arabic*.]
  \item En la Parte~A hay una sola opción correcta por cuestión. Las respuestas erróneas no restan.
  \item En la Parte~B deben justificarse todos los pasos y escribirse las unidades. Los resultados finales se dan con tres cifras significativas y los intermedios con cuatro. Cuando haga falta $\Phi$ o $\varphi$ fuera de los valores tabulados se interpola linealmente; se admiten diferencias de $\pm0{,}002$ debidas a la interpolación.
  \item Todos los datos de empresas, plantas, escenas y flotas son sintéticos. Las cifras de casos industriales publicados que aparecen en los apuntes no se preguntan de memoria.
  \item Salvo indicación contraria, los problemas de optimización bayesiana son de \emph{minimización} y $\fmejor$ es el menor valor observado, como en M25, tema~4.
\end{enumerate}

\section*{Datos de referencia}

\begin{table}[h]
\centering\footnotesize
\setlength{\tabcolsep}{4pt}
\caption{Densidad $\varphi(z)$ y función de distribución $\Phi(z)$ de la normal estándar. Para $z<0$: $\varphi(-z)=\varphi(z)$ y $\Phi(-z)=1-\Phi(z)$.}\label{tab:normal}
\begin{tabular}{@{}ccc@{\qquad}ccc@{\qquad}ccc@{}}
\toprule
$z$ & $\varphi(z)$ & $\Phi(z)$ & $z$ & $\varphi(z)$ & $\Phi(z)$ & $z$ & $\varphi(z)$ & $\Phi(z)$ \\
\midrule
0{,}00 & 0{,}3989 & 0{,}5000 & 0{,}85 & 0{,}2780 & 0{,}8023 & 1{,}70 & 0{,}0940 & 0{,}9554 \\
0{,}05 & 0{,}3984 & 0{,}5199 & 0{,}90 & 0{,}2661 & 0{,}8159 & 1{,}75 & 0{,}0863 & 0{,}9599 \\
0{,}10 & 0{,}3970 & 0{,}5398 & 0{,}95 & 0{,}2541 & 0{,}8289 & 1{,}80 & 0{,}0790 & 0{,}9641 \\
0{,}15 & 0{,}3945 & 0{,}5596 & 1{,}00 & 0{,}2420 & 0{,}8413 & 1{,}85 & 0{,}0721 & 0{,}9678 \\
0{,}20 & 0{,}3910 & 0{,}5793 & 1{,}05 & 0{,}2299 & 0{,}8531 & 1{,}90 & 0{,}0656 & 0{,}9713 \\
0{,}25 & 0{,}3867 & 0{,}5987 & 1{,}10 & 0{,}2179 & 0{,}8643 & 1{,}95 & 0{,}0596 & 0{,}9744 \\
0{,}30 & 0{,}3814 & 0{,}6179 & 1{,}15 & 0{,}2059 & 0{,}8749 & 2{,}00 & 0{,}0540 & 0{,}9772 \\
0{,}35 & 0{,}3752 & 0{,}6368 & 1{,}20 & 0{,}1942 & 0{,}8849 & 2{,}05 & 0{,}0488 & 0{,}9798 \\
0{,}40 & 0{,}3683 & 0{,}6554 & 1{,}25 & 0{,}1826 & 0{,}8944 & 2{,}10 & 0{,}0440 & 0{,}9821 \\
0{,}45 & 0{,}3605 & 0{,}6736 & 1{,}30 & 0{,}1714 & 0{,}9032 & 2{,}15 & 0{,}0396 & 0{,}9842 \\
0{,}50 & 0{,}3521 & 0{,}6915 & 1{,}35 & 0{,}1604 & 0{,}9115 & 2{,}20 & 0{,}0355 & 0{,}9861 \\
0{,}55 & 0{,}3429 & 0{,}7088 & 1{,}40 & 0{,}1497 & 0{,}9192 & 2{,}25 & 0{,}0317 & 0{,}9878 \\
0{,}60 & 0{,}3332 & 0{,}7257 & 1{,}45 & 0{,}1394 & 0{,}9265 & 2{,}30 & 0{,}0283 & 0{,}9893 \\
0{,}65 & 0{,}3230 & 0{,}7422 & 1{,}50 & 0{,}1295 & 0{,}9332 & 2{,}35 & 0{,}0252 & 0{,}9906 \\
0{,}70 & 0{,}3123 & 0{,}7580 & 1{,}55 & 0{,}1200 & 0{,}9394 & 2{,}40 & 0{,}0224 & 0{,}9918 \\
0{,}75 & 0{,}3011 & 0{,}7734 & 1{,}60 & 0{,}1109 & 0{,}9452 & 2{,}45 & 0{,}0198 & 0{,}9929 \\
0{,}80 & 0{,}2897 & 0{,}7881 & 1{,}65 & 0{,}1023 & 0{,}9505 & 2{,}50 & 0{,}0175 & 0{,}9938 \\
\bottomrule
\end{tabular}
\end{table}

\begin{table}[h]
\centering\small
\caption{Cuantiles y constantes.}\label{tab:cuantiles}
\begin{tabular}{@{}lc@{\qquad}lc@{}}
\toprule
Cantidad & Valor & Cantidad & Valor \\
\midrule
$z_{0{,}75}=\Phi^{-1}(0{,}75)$ & $0{,}6745$ & $t_{15;\,0{,}975}$ (Student, 15 g.\,l.) & $2{,}131$ \\
$z_{0{,}80}$ & $0{,}8416$ & $t_{15;\,0{,}95}$ & $1{,}753$ \\
$z_{0{,}95}$ & $1{,}645$ & $e^{-1/2}$ & $0{,}6065$ \\
$z_{0{,}975}$ & $1{,}960$ & $e^{-2}$ & $0{,}1353$ \\
$\ln 0{,}95$ & $-0{,}05129$ & $e^{-9/2}$ & $0{,}01111$ \\
\bottomrule
\end{tabular}
\end{table}

\section*{Fórmulas}
\[
\EI_n(x)=(\fmejor-\mu)\,\Phi(z)+\sigma\,\varphi(z),\quad z=\frac{\fmejor-\mu}{\sigma};
\]
\[
\mu_n(x)=k_\ast\T K_y^{-1}y,\qquad \sigma_n^2(x)=k(x,x)-k_\ast\T K_y^{-1}k_\ast,\qquad K_y=K+\sigma_\varepsilon^{2}I ;
\]
\[
\VAN=-I+b\,\frac{1-(1+\rho)^{-N}}{\rho};\qquad
n\approx\Bigl(\frac{(z_{1-\alpha/2}+z_{1-\beta})\,\sigma_D}{\delta}\Bigr)^{2};\qquad
\bar d\pm t_{n-1,\,1-\alpha/2}\,\frac{s_d}{\sqrt n}.
\]
